\subsection{Potential Energy}
The potential energy is the work done in moving the particle from a radius of $\infty$ to $r$.
So the work done is given by
\begin{align}
U(r)&=\int_{\infty}^{r} \bm{F}(r') \cdot d\bm{r}'  \nonumber \\
U(r)&=\int_{\infty}^{r} -\frac{GMm}{r'^2} \hat{\bm{r}} \cdot dr' \hat{\bm{r}}  \nonumber \\
U(r)&=-GMm\int_{\infty}^{r} \frac{dr'}{r'^2} \nonumber \\
U(r)&=-GMm \left[ \frac{-1}{r'} \right]_{\infty}^{r} \nonumber \\
U(r)&=GMm \left[ \frac{1}{\infty} - \frac{1}{r} \right] \nonumber \\
U(r)&=- \frac{GMm}{r} \label{eq:potentialEnergy}
\end{align}

\subsection{Kinetic Energy}
The kinetic energy of a particle is given by
\begin{equation}
T(\dot{r})=\frac{1}{2}m\dot{\bm{r}}\cdot\dot{\bm{r}} \label{eq:KineticEnergy}
\end{equation}
using eqation \ref{eq:vel} gives 
\begin{align}
T(\dot{r})&=\frac{1}{2}m\left[ \dot{r}\hat{\bm{r}}+r\omega\hat{\bm{\theta}}\right]\cdot\left[ \dot{r}\hat{\bm{r}}+r\omega\hat{\bm{\theta}}\right] \nonumber \\
&=\frac{1}{2}m \dot{r}^2 + \frac{1}{2}mr^2 \omega^2 \label{eq:Gravity:KineticEnergy}
\end{align}

we can also make the subsitution using using the magnitude of the angular momentum\footnote{equation \ref{eq:AngularMomentumScalar}} $l=mr^2\omega$

\begin{align}
l&=mr^2\omega \nonumber \\
\omega&=\frac{l}{mr^2} \label{eq:Gravity:AngularMomentum}
\end{align}

using equation \ref{eq:Gravity:AngularMomentum} in equation \ref{eq:Gravity:KineticEnergy} leads to

\begin{align}
T(\dot{r})&=\frac{1}{2}m \dot{r}^2 + \frac{1}{2}mr^2 \omega^2 \nonumber \\
&=\frac{1}{2}m \dot{r}^2 + \frac{1}{2}mr^2 \frac{l^2}{m^2r^4} \nonumber \\
&=\frac{1}{2}m \dot{r}^2 + \frac{l^2}{2mr^2} \label{eq:Gravity:KineticEnergyFinal}
\end{align}

\subsubsection{Total Energy}
Finally the total energy can be found by putting together the equation for kinetic energy and potential energy
\begin{align}
E(r,\dot{r})=T(\dot{r})+U(r) \nonumber \\
E(r,\dot{r})=\frac{1}{2}m \dot{r}^2 + \frac{l^2}{2mr^2}- \frac{GMm}{r}\label{eq:Gravity:TotalEnergy}
\end{align}

sometimes the last two terms of equation \ref{eq:Gravity:TotalEnergy} are put together to form a new potential energy type term that is dependent only on the radius and the angular momentum.

\subsection{Conservation of Total Energy}

Just like with angular momentum we can show that total energy is conserved by taking the derivative with respect to time.

\begin{align}
\frac{dE}{dt}&=\frac{d}{dt}\left[\frac{1}{2}m \dot{r}^2 +  \frac{l^2}{2mr^2}-\frac{GMm}{r}\right] \nonumber \\
&=\frac{1}{2}m \frac{d\dot{r}^2}{dt} +  \frac{l^2}{2m}\frac{dr^{-2}}{dt}-GMm\frac{dr^{-1}}{dt} \nonumber \\
&=\frac{1}{2}m \frac{d\dot{r}^2}{d\dot{r}}\frac{d\dot{r}}{dt} +  \frac{l^2}{2m}\frac{dr^{-2}}{dr}\frac{dr}{dt}-GMm\frac{dr^{-1}}{dr}\frac{dr}{dt} \nonumber \\
&=\frac{1}{2}m (2\dot{r})\ddot{r}+ \frac{l^2}{2m}\left(\frac{-2}{r^3}\right)\dot{r}+\frac{GMm}{r^2}\dot{r} \nonumber \\
&=m \dot{r}\ddot{r}- \frac{l^2}{mr^3}\dot{r}+\frac{GMm}{r^2}\dot{r}
\end{align}

using the magnitude of the angular momentum\footnote{equation \ref{eq:AngularMomentumScalar}} $l=mr^2\omega$

\begin{align}
\frac{dE}{dt} &=m \dot{r}\ddot{r}- \frac{l^2}{mr^3}\dot{r}+\frac{GMm}{r^2}\dot{r} \nonumber \\
&=m \dot{r}\ddot{r}- \frac{m^2r^4\omega^2}{mr^3}\dot{r}+\frac{GMm}{r^2}\dot{r} \nonumber \\
&=m \dot{r}\ddot{r}-mr\omega^2\dot{r}+\frac{GMm}{r^2}\dot{r} \nonumber \\
&=m \dot{r}\left[\ddot{r}-r\omega^2+\frac{GM}{r^2}\right]
\end{align}

using the radial part of the acceleration\footnote{equation \ref{eq:radialAcceleration}} $-\frac{GM}{r^2}= \ddot{r} -r \omega^2$ we get
\begin{align}
\frac{dE}{dt}&=m \dot{r}\left[-\frac{GM}{r^2}+\frac{GM}{r^2}\right] \nonumber \\
&=m \dot{r}\left[0\right] \nonumber \\
&=0
\end{align}

This then shows that the total energy does not change with respect to time, so is a constant and also conserved.